\begin{code}[hide]%
\>[0]\AgdaKeyword{module}\AgdaSpace{}%
\AgdaModule{paper.sections.examples}\AgdaSpace{}%
\AgdaKeyword{where}\<%
\\
%
\\[\AgdaEmptyExtraSkip]%
\>[0]\AgdaKeyword{open}\AgdaSpace{}%
\AgdaKeyword{import}\AgdaSpace{}%
\AgdaModule{Algebra.Bundles}\AgdaSpace{}%
\AgdaKeyword{using}\AgdaSpace{}%
\AgdaSymbol{(}\AgdaRecord{Group}\AgdaSpace{}%
\AgdaSymbol{;}\AgdaSpace{}%
\AgdaRecord{Monoid}\AgdaSymbol{)}\<%
\\
\>[0]\AgdaKeyword{open}\AgdaSpace{}%
\AgdaKeyword{import}\AgdaSpace{}%
\AgdaModule{Algebra.Morphism.Structures}\AgdaSpace{}%
\AgdaKeyword{using}\AgdaSpace{}%
\AgdaSymbol{(}\AgdaKeyword{module}\AgdaSpace{}%
\AgdaModule{MonoidMorphisms}\AgdaSymbol{)}\<%
\\
\>[0]\AgdaKeyword{open}\AgdaSpace{}%
\AgdaKeyword{import}\AgdaSpace{}%
\AgdaModule{Data.Nat}\AgdaSpace{}%
\AgdaKeyword{using}\AgdaSpace{}%
\AgdaSymbol{(}\AgdaDatatype{ℕ}\AgdaSpace{}%
\AgdaSymbol{;}\AgdaSpace{}%
\AgdaInductiveConstructor{suc}\AgdaSpace{}%
\AgdaSymbol{;}\AgdaSpace{}%
\AgdaInductiveConstructor{2+}\AgdaSymbol{)}\<%
\\
\>[0]\AgdaKeyword{open}\AgdaSpace{}%
\AgdaKeyword{import}\AgdaSpace{}%
\AgdaModule{Data.Nat.Primality}\AgdaSpace{}%
\AgdaKeyword{using}\AgdaSpace{}%
\AgdaSymbol{(}\AgdaRecord{Prime}\AgdaSymbol{)}\<%
\\
\>[0]\AgdaKeyword{open}\AgdaSpace{}%
\AgdaKeyword{import}\AgdaSpace{}%
\AgdaModule{Level}\AgdaSpace{}%
\AgdaKeyword{using}\AgdaSpace{}%
\AgdaSymbol{(}\AgdaFunction{0ℓ}\AgdaSymbol{)}\<%
\\
%
\\[\AgdaEmptyExtraSkip]%
\>[0]\AgdaKeyword{open}\AgdaSpace{}%
\AgdaKeyword{import}\AgdaSpace{}%
\AgdaModule{Notations}\AgdaSpace{}%
\AgdaKeyword{using}\AgdaSpace{}%
\AgdaSymbol{(}\AgdaInductiveConstructor{₁₊}\AgdaSymbol{)}\<%
\\
\>[0]\AgdaKeyword{open}\AgdaSpace{}%
\AgdaKeyword{import}\AgdaSpace{}%
\AgdaModule{Word.Base}\AgdaSpace{}%
\AgdaKeyword{using}\AgdaSpace{}%
\AgdaSymbol{(}\AgdaOperator{\AgdaFunction{\AgdaUnderscore{}ʷ}}\AgdaSymbol{)}\<%
\\
\>[0]\AgdaKeyword{open}\AgdaSpace{}%
\AgdaKeyword{import}\AgdaSpace{}%
\AgdaModule{Presentation.Definitions}\AgdaSpace{}%
\AgdaKeyword{using}\AgdaSpace{}%
\AgdaSymbol{(}\AgdaOperator{\AgdaRecord{\AgdaUnderscore{}IsPresentationOf\AgdaUnderscore{}}}\AgdaSymbol{)}\<%
\\
\>[0]\AgdaKeyword{open}\AgdaSpace{}%
\AgdaKeyword{import}\AgdaSpace{}%
\AgdaModule{Presentation.Construct.Base}\AgdaSpace{}%
\AgdaKeyword{using}\AgdaSpace{}%
\AgdaSymbol{(}\AgdaOperator{\AgdaDatatype{\AgdaUnderscore{}⋄\AgdaUnderscore{}⋄\AgdaUnderscore{}}}\AgdaSpace{}%
\AgdaSymbol{;}\AgdaSpace{}%
\AgdaDatatype{CommRel}\AgdaSpace{}%
\AgdaSymbol{;}\AgdaSpace{}%
\AgdaDatatype{ConjRelʷ}\AgdaSpace{}%
\AgdaSymbol{;}\AgdaSpace{}%
\AgdaOperator{\AgdaFunction{\AgdaUnderscore{}⊕\textasciicircum{}\AgdaUnderscore{}}}\AgdaSymbol{)}\<%
\\
\>[0]\AgdaKeyword{import}\AgdaSpace{}%
\AgdaModule{Presentation.Properties}\AgdaSpace{}%
\AgdaSymbol{as}\AgdaSpace{}%
\AgdaModule{PP}\<%
\\
\>[0]\AgdaKeyword{import}\AgdaSpace{}%
\AgdaModule{Presentation.Construct.Properties.DirectProduct}\AgdaSpace{}%
\AgdaSymbol{as}\AgdaSpace{}%
\AgdaModule{DP}\<%
\\
\>[0]\AgdaKeyword{import}\AgdaSpace{}%
\AgdaModule{Presentation.Construct.Properties.NDirectProduct}\AgdaSpace{}%
\AgdaSymbol{as}\AgdaSpace{}%
\AgdaModule{NDP}\<%
\\
\>[0]\AgdaKeyword{open}\AgdaSpace{}%
\AgdaKeyword{import}\AgdaSpace{}%
\AgdaModule{Examples.Groups.Cyclic.Normalization}\AgdaSpace{}%
\AgdaKeyword{using}\AgdaSpace{}%
\AgdaSymbol{(}\AgdaOperator{\AgdaDatatype{\AgdaUnderscore{}Cn,\AgdaUnderscore{}===\AgdaUnderscore{}}}\AgdaSymbol{)}\<%
\\
\>[0]\AgdaKeyword{open}\AgdaSpace{}%
\AgdaKeyword{import}\AgdaSpace{}%
\AgdaModule{Examples.Groups.Cyclic.Semantics}\AgdaSpace{}%
\AgdaKeyword{using}\AgdaSpace{}%
\AgdaSymbol{(}\AgdaFunction{Cn-group}\AgdaSymbol{)}\<%
\\
\>[0]\AgdaKeyword{import}\AgdaSpace{}%
\AgdaModule{Examples.Groups.Cyclic.Presentation}\AgdaSpace{}%
\AgdaSymbol{as}\AgdaSpace{}%
\AgdaModule{CyP}\<%
\\
\>[0]\AgdaKeyword{open}\AgdaSpace{}%
\AgdaKeyword{import}\AgdaSpace{}%
\AgdaModule{Examples.Groups.Symmetric.Syntactics}\AgdaSpace{}%
\AgdaKeyword{using}\AgdaSpace{}%
\AgdaSymbol{(}\AgdaOperator{\AgdaDatatype{\AgdaUnderscore{}VRel,\AgdaUnderscore{}===\AgdaUnderscore{}}}\AgdaSymbol{)}\<%
\\
\>[0]\AgdaKeyword{import}\AgdaSpace{}%
\AgdaModule{Examples.Groups.Pauli.Presentation}\AgdaSpace{}%
\AgdaSymbol{as}\AgdaSpace{}%
\AgdaModule{Pauli}\<%
\\
\>[0]\AgdaKeyword{import}\AgdaSpace{}%
\AgdaModule{Examples.Construct.SemiDirectProduct.SnD}\AgdaSpace{}%
\AgdaSymbol{as}\AgdaSpace{}%
\AgdaModule{SnD}\<%
\\
\>[0]\AgdaKeyword{import}\AgdaSpace{}%
\AgdaModule{Examples.Amalgamations.CliffordT1}\AgdaSpace{}%
\AgdaSymbol{as}\AgdaSpace{}%
\AgdaModule{CliffordT1}\<%
\\
%
\\[\AgdaEmptyExtraSkip]%
\>[0]\AgdaKeyword{open}\AgdaSpace{}%
\AgdaModule{SnD}\AgdaSpace{}%
\AgdaKeyword{using}\AgdaSpace{}%
\AgdaSymbol{(}\AgdaFunction{conj}\AgdaSpace{}%
\AgdaSymbol{;}\AgdaSpace{}%
\AgdaKeyword{module}\AgdaSpace{}%
\AgdaModule{Wreath}\AgdaSymbol{)}\<%
\\
\>[0]\AgdaKeyword{open}\AgdaSpace{}%
\AgdaModule{Monoid}\AgdaSpace{}%
\AgdaKeyword{using}\AgdaSpace{}%
\AgdaSymbol{(}\AgdaFunction{rawMonoid}\AgdaSymbol{)}\<%
\\
\>[0]\AgdaKeyword{open}\AgdaSpace{}%
\AgdaModule{MonoidMorphisms}\AgdaSpace{}%
\AgdaKeyword{using}\AgdaSpace{}%
\AgdaSymbol{(}\AgdaRecord{IsMonoidIsomorphism}\AgdaSymbol{)}\<%
\\
\>[0]\AgdaKeyword{open}\AgdaSpace{}%
\AgdaModule{PP}\AgdaSpace{}%
\AgdaKeyword{using}\AgdaSpace{}%
\AgdaSymbol{(}\AgdaFunction{•-ε-monoid}\AgdaSymbol{)}\<%
\\
\>[0]\AgdaKeyword{open}\AgdaSpace{}%
\AgdaModule{CliffordT1.CliffordT1}\AgdaSpace{}%
\AgdaKeyword{using}\AgdaSpace{}%
\AgdaSymbol{(}\AgdaOperator{\AgdaDatatype{\AgdaUnderscore{}===\AgdaUnderscore{}}}\AgdaSpace{}%
\AgdaSymbol{;}\AgdaSpace{}%
\AgdaFunction{amalPres}\AgdaSpace{}%
\AgdaSymbol{;}\AgdaSpace{}%
\AgdaFunction{f}\AgdaSymbol{)}\<%
\end{code}
\section{Examples: A Catalogue of Verified Presentations}\label{sec:examples}

This section is the library's harvest, but the inventory is not the point;
what each entry was chosen to exercise is.  The families below were
selected so that every layer of \S\ref{sec:design} is load-bearing
somewhere.  Cyclic groups are the base case of every tower and the
smallest place where monoid and group presentations come apart.  Symmetric
groups exercise the circuit layer and a genuinely recursive coset tower.
The product demonstration and the wreath products exercise the direct and
semidirect constructions with their semantic theorems, and yield a family
of completeness results parameterized over two naturals.  The Pauli groups
are the pure-composition test: a presentation with no coset enumeration
at all.  The three amalgamation studies --- qubit Clifford+$T$, qutrit
Clifford+$T$, and $U_3(\mathbb{Z}[\tfrac12,i])$ --- are the stress tests:
longer towers, twisted relations, branching chains, definitional
extensions, and the largest \aid{refl}-discharged obligation sets in the
development.  For each family we say what its normal form \emph{is},
since that is what a reader who wants to reuse the machinery on a new
gate set needs to see first.

Table~\ref{tab:presentations} lists every concrete
\aid{{\_}IsPresentationOf{\_}} theorem in the development;
Table~\ref{tab:nfs} lists the normal-form and uniqueness witnesses behind
them.  All of them are restated with full types in the index module
\texttt{MainTheorems.agda}, which also re-exports the conditional
construction theorems of \S\ref{sec:constructions} with their premise
telescopes.  We then discuss each family, ending with the three
amalgamation case studies, whose final theorems are monoid
\emph{isomorphisms} between gate-set presentations and amalgamated
products.

\begin{table}
\caption{Concrete presentation theorems (all verified; names as in
\texttt{MainTheorems.agda}).  $p$ ranges over odd primes, $n, m, N$ over
naturals.}
\label{tab:presentations}
\centering
\small
\begin{tabular}{@{}lll@{}}
\toprule
Relation set & Presents & Index name \\
\midrule
\aid{EmptyRel} over $\bot$ & trivial group & \aid{trivial-presentation} \\
\aid{TrivialRel} over any $A$ & trivial group & \aid{trivial-presentation′} \\
$T^{\,n+1} = \varepsilon$ & $\mathbb{Z}/(n{+}1)\mathbb{Z}$ & \aid{cyclic-presentation} \\
$T^{\,0} = \varepsilon$ & the \emph{monoid} $(\mathbb{N}, +, 0)$ & \aid{cyclic-monoid-presentation} \\
\aid{order}, \aid{yang-baxter} + structural & permutations of $\mathsf{Fin}\;n$ & \aid{symmetric-presentation} \\
$C_p \aidop{⋄} C_p \aidop{⋄} \aid{CommRel}$ & $\mathbb{Z}_p \times \mathbb{Z}_p$ & (\aid{H-pres}, Pauli development) \\
$(C_p \aidop{⋄} C_p \aidop{⋄} \aid{CommRel}) \aidop{⊕\char94} n$ & Pauli quotient $(\mathbb{Z}_p{\times}\mathbb{Z}_p)^n$ & \aid{pauli-presentation} \\
$(C_N \aidop{⊕\char94} n) \aidop{⋄} S_n\text{-circuits} \aidop{⋄} \aid{ConjRelʷ}$ & $\mathbb{Z}_N \wr S_n$ (signed permutations) & \aid{wreath-presentation} \\
\bottomrule
\end{tabular}
\end{table}

\begin{table}
\caption{Normal-form and uniqueness witnesses (selection).  NFI =
\aid{NormalFormInjective}, NF = \aid{NormalForm} (with section), UNF =
\aid{UniqueNormalForm}.  For the three constructions, UNF is lifted
from the factor witnesses --- for the amalgam, from two first-order
exactness certificates (base normal form, coset table on section
words).}
\label{tab:nfs}
\centering
\footnotesize
\setlength{\tabcolsep}{4pt}
\begin{tabular}{@{}llll@{}}
\toprule
Presentation & Carrier & Witnesses & Note \\
\midrule
cyclic, order $n{+}1$ & $\mathsf{Fin}\,(n{+}1)$ & NFI, NF, UNF & presents $\mathbb{Z}/(n{+}1)\mathbb{Z}$ \\
cyclic, order $0$ & $\mathbb{N}$ & NFI, NF, UNF & presents the monoid $(\mathbb{N}, +, 0)$ (\S\ref{sec:cyclic}) \\
$S_n$ circuits & $\top \times C\,1 \times \dots \times C\,(n{-}1)$ & NFI, NF, UNF$\times 2$ & factorial system, $|{\cdot}| = n!$ \\
$\Gamma \oplus \Delta$ & $\mathit{NF}_1 \times \mathit{NF}_2$ & NFI, NF, UNF & lifted from factors \\
$\Gamma \rtimes \Delta \star \mathit{conj}$ & $\mathit{NF}_1 \times \mathit{NF}_2$ & NFI, NF, UNF & twist on the left factor \\
amalgamated $P_1 \ast P_2$ & $(D \uplus \top) \times \mathsf{List}(C {\times} D) \times (C \uplus \top)$ & NFI, NF, UNF & alternating cosets \\
Clifford+$T$ base $\langle \omega \rangle {\times} \langle S \rangle$ & $\mathsf{Fin}\,8 \times \mathsf{Fin}\,4$ & NFI, NF, UNF & radices $8, 4$; unique for $\mathbb{Z}_8 {\times} \mathbb{Z}_4$ \\
\bottomrule
\end{tabular}
\end{table}

\subsection{Cyclic groups, and the trivial group twice}\label{sec:cyclic}

The cyclic development presents $\langle\, T \mid T^{\,n+1} =
\varepsilon\,\rangle$ with normal-form carrier $\mathsf{Fin}\,(n{+}1)$,
for every positive order at once (\aid{cyclic-presentation}).  The normal
form of a word is its letter count reduced modulo $n{+}1$, and the section
sends the residue $k$ back to the word $T^{k}$; uniqueness for the
semantics $\mathbb{Z}/(n{+}1)\mathbb{Z}$ is immediate.  Order
$0$ gets a theorem of its own, because it is the smallest example
separating \emph{monoid} presentations from \emph{group} presentations
--- and every presentation in this library is a monoid presentation
(\S\ref{sec:presentations}).  Read as a group presentation, the
trivial relation $T^{\,0} = \varepsilon$ presents the cyclic group of
order $0$, which is by definition the infinite cyclic group
$\mathbb{Z}$.  The monoid it presents, however, is the free monoid on
one generator, and that is what the library proves:
\aid{cyclic-monoid-presentation} states that $T^{\,0} = \varepsilon$
\aid{IsMonoidPresentationOf} $(\mathbb{N}, +, 0)$, with the same
uniform normal-form witness (carrier $\mathbb{N}$, a word count with
no wraparound).  The gap between the two readings is precisely
\aid{Grouplike}: a grouplike monoid presentation \emph{is} a group
presentation (\aid{monoid-presentation⇒group-presentation}), and at
order $0$ the generator has no inverse in the presented monoid, so the
upgrade --- available at every positive order --- is exactly what
fails.  The
trivial group is presented twice --- empty alphabet with no relations, and
arbitrary alphabet with the everything-is-$\varepsilon$ relation --- and
the two presentations are proved isomorphic; besides being the base case
of $n$-fold products, this pair is a two-page worked example of the whole
pipeline and a good entry point into the code.

\subsection{A product demonstration: $S_3 \times C_5$}

A small worked example composes two unrelated developments: the direct
product of the $3$-wire permutation presentation and the cyclic
presentation of order $5$, with the lifted normal form of
\S\ref{sec:constructions} on the carrier $\mathit{NF}_{S_3} \times
\mathsf{Fin}\,5$ --- a normal form is a pair of a factorial digit tuple
and a residue, and its section is the concatenation of the two factor
sections.  The example exists to be read: it shows the transport
machinery end to end on a group of order $30$, with the monoid-level
semantics assembled from the factor semantics.

\subsection{Wreath products: signed permutations and beyond}

The headline compositional result combines almost every construction in
the library.  Fix a base order $N = m{+}1$ and $n$ wires.  The $n$-fold
direct power $(C_N)^{\oplus n}$ presents $(\mathbb{Z}_N)^n$; the circuit
presentation of $S_n$ acts on it by permuting coordinates --- an action
specified \emph{syntactically} as a word-valued conjugation
$\mathit{conj} : \aid{Gen}\;n \to (\top \uplus^n) \to
\aid{Word}\;(\top \uplus^n)$ sending a generator and a coordinate to the
permuted coordinate.  Two finite compatibility lemmas (the action
respects the Coxeter relations pointwise, and respects the product
relations) feed the semidirect-product presentation theorem, giving
\begin{code}%
\>[0]\AgdaFunction{wreath-presentation}\AgdaSpace{}%
\AgdaSymbol{:}\AgdaSpace{}%
\AgdaSymbol{∀}\AgdaSpace{}%
\AgdaBound{n}\AgdaSpace{}%
\AgdaBound{m}\AgdaSpace{}%
\AgdaSymbol{→}%
\>[111I]\AgdaSymbol{(((}\AgdaInductiveConstructor{suc}\AgdaSpace{}%
\AgdaBound{m}\AgdaSpace{}%
\AgdaOperator{\AgdaDatatype{Cn,\AgdaUnderscore{}===\AgdaUnderscore{}}}\AgdaSymbol{)}\AgdaSpace{}%
\AgdaOperator{\AgdaFunction{⊕\textasciicircum{}}}\AgdaSpace{}%
\AgdaBound{n}\AgdaSymbol{)}\AgdaSpace{}%
\AgdaOperator{\AgdaDatatype{⋄}}\AgdaSpace{}%
\AgdaSymbol{(}\AgdaBound{n}\AgdaSpace{}%
\AgdaOperator{\AgdaDatatype{VRel,\AgdaUnderscore{}===\AgdaUnderscore{}}}\AgdaSymbol{)}\AgdaSpace{}%
\AgdaOperator{\AgdaDatatype{⋄}}\AgdaSpace{}%
\AgdaDatatype{ConjRelʷ}\AgdaSpace{}%
\AgdaFunction{conj}\AgdaSymbol{)}\<%
\\
\>[111I][@{}l@{\AgdaIndent{0}}]%
\>[32]\AgdaOperator{\AgdaRecord{IsPresentationOf}}\AgdaSpace{}%
\AgdaSymbol{(}\AgdaFunction{Wreath.wreath-group}\AgdaSpace{}%
\AgdaBound{n}\AgdaSpace{}%
\AgdaBound{m}\AgdaSymbol{)}\<%
\end{code}
\begin{code}[hide]%
\>[0]\AgdaFunction{wreath-presentation}\AgdaSpace{}%
\AgdaBound{n}\AgdaSpace{}%
\AgdaBound{m}\AgdaSpace{}%
\AgdaSymbol{=}\AgdaSpace{}%
\AgdaFunction{Wreath.presentation}\AgdaSpace{}%
\AgdaBound{n}\AgdaSpace{}%
\AgdaBound{m}\<%
\end{code}
where $\aid{wreath-group}\;n\;m$ is the semantic
$(\mathbb{Z}_N)^n \rtimes S_n = \mathbb{Z}_N \wr S_n$ built from the
\texttt{ForStdlib} semidirect product and the bundled permutation group.
The lifted normal form is a pair: $n$ residues modulo $N$, one phase digit
per wire, followed by the factorial tuple of a permutation --- so its
section is a word of $T$-powers, one per coordinate, followed by the
staircases of \S\ref{sec:permutations}, and the coset digits of the two
factors simply concatenate, with the twist of the semidirect product
accounted for in the transported table.  At $N = 2$ this is the group of
\emph{signed permutations} --- the hyperoctahedral group $B_n$, the
symmetry group of the hypercube; at general $N$ it is the generalized
symmetric group $G(N,1,n)$ of monomial matrices with $N$-th-root-of-unity
entries.  These groups are ubiquitous in the qudit setting (Pauli-frame
changes, permutation-and-phase subgroups of Clifford groups), and the
theorem delivers their presentations \emph{for all $N$ and $n$ at once}
--- a family of completeness results parameterized over two naturals,
which is where a mechanized, compositional treatment clearly outruns the
per-instance pen-and-paper approach.

\subsection{Pauli groups in all odd prime dimensions}\label{sec:pauli}

The Pauli development is the purest showcase of compositionality --- there
is \emph{no} coset enumeration in it at all:

\begin{code}%
\>[0]\AgdaFunction{pauli-presentation}\AgdaSpace{}%
\AgdaSymbol{:}\AgdaSpace{}%
\AgdaSymbol{∀}\AgdaSpace{}%
\AgdaSymbol{(}\AgdaBound{p-2}\AgdaSpace{}%
\AgdaSymbol{:}\AgdaSpace{}%
\AgdaDatatype{ℕ}\AgdaSymbol{)}\AgdaSpace{}%
\AgdaSymbol{(}\AgdaBound{p-prime}\AgdaSpace{}%
\AgdaSymbol{:}\AgdaSpace{}%
\AgdaRecord{Prime}\AgdaSpace{}%
\AgdaSymbol{(}\AgdaInductiveConstructor{2+}\AgdaSpace{}%
\AgdaBound{p-2}\AgdaSymbol{))}\AgdaSpace{}%
\AgdaBound{n}\AgdaSpace{}%
\AgdaSymbol{→}\<%
\\
\>[0][@{}l@{\AgdaIndent{0}}]%
\>[2]\AgdaSymbol{(}\AgdaFunction{Pauli.Γ-H}\AgdaSpace{}%
\AgdaBound{p-2}\AgdaSpace{}%
\AgdaBound{p-prime}\AgdaSpace{}%
\AgdaOperator{\AgdaFunction{⊕\textasciicircum{}}}\AgdaSpace{}%
\AgdaBound{n}\AgdaSymbol{)}\AgdaSpace{}%
\AgdaOperator{\AgdaRecord{IsPresentationOf}}\AgdaSpace{}%
\AgdaSymbol{(}\AgdaFunction{Pauli.Pauli-group}\AgdaSpace{}%
\AgdaBound{p-2}\AgdaSpace{}%
\AgdaBound{p-prime}\AgdaSpace{}%
\AgdaBound{n}\AgdaSymbol{)}\<%
\end{code}
\begin{code}[hide]%
\>[0]\AgdaFunction{pauli-presentation}\AgdaSpace{}%
\AgdaBound{p-2}\AgdaSpace{}%
\AgdaBound{p-prime}\AgdaSpace{}%
\AgdaSymbol{=}\AgdaSpace{}%
\AgdaFunction{Pauli.Pauli-presentation}\AgdaSpace{}%
\AgdaBound{p-2}\AgdaSpace{}%
\AgdaBound{p-prime}\<%
\end{code}

For an odd prime $p$ --- a \emph{qupit} is a qudit whose dimension is an
odd prime $p$ --- the single-qupit relation set is the direct-product
join of two order-$p$ cyclic presentations ($X$- and $Z$-type generators,
which commute in the quotient of the Pauli group by its phase subgroup),
and the $n$-qupit relation set is its $n$-fold power.  The semantic group
$(\mathbb{Z}_p \times \mathbb{Z}_p)^n$ is assembled by the standard
library's binary direct product and our $n$-fold iteration, and the
presentation is literally three construction theorems composed --- the
cyclic presentation, the binary-product theorem, and the $n$-fold-product
theorem, in three lines of Agda.  The normal form is correspondingly a
tuple of $n$ pairs of exponents
$(a_i, b_i) \in \mathsf{Fin}\,p \times \mathsf{Fin}\,p$, one per wire,
with section $X^{a_1} Z^{b_1} \cdots X^{a_n} Z^{b_n}$; nothing about it
was proved for the Pauli group specifically.  The full Pauli group with
phases, and its role beneath the Clifford hierarchy, is left to future
work (\S\ref{sec:conclusion}).

\subsection{Single-qubit Clifford+$T$, as an amalgamated free product}\label{sec:cliffordt1}

The first amalgamation case study presents the single-qubit Clifford+$T$
\emph{monoid}, the flagship example of the completeness literature for
good reasons: Clifford+$T$ is the smallest universal gate set in common
use, its single-qubit fragment has the first and best-known canonical
form --- the Matsumoto--Amano normal
form~\cite{matsumoto2008representation,giles2013remarks}, which underlies
exact synthesis and $T$-count optimality --- and every later normal form
for Clifford-style gate sets is modelled on it.  The relation set, over
generators $T, H, S, \omega$ with $X := H S S H$ as an abbreviation, is
\[
  \omega^8 = \varepsilon,\quad
  T^2 = S,\quad
  S^4 = \varepsilon,\quad
  H^2 = \varepsilon,\quad
  (SH)^3 = \omega,\quad
  (TX)^2 = \omega,\quad
  \omega\ \text{central},
\]
and the verified structure exhibits the monoid as an amalgamated free
product over the subgroup $M$ generated by $X$, $S$, and $\omega$:
\[
  \mathrm{Clifford}{+}T \;\cong\;
  \langle T, M \rangle \,\ast_{M}\, \langle H, M \rangle,
  \qquad M = \langle X, S, \omega \rangle .
\]
The tower starts from the direct product $\langle \omega \rangle \times
\langle S \rangle$ of cyclic groups of orders $8$ and $4$ (normal-form
radices $8 \cdot 4$), extends it by a two-coset table ($\varepsilon$,
$X$) to $M$, and then extends \emph{twice from $M$}: by a two-coset
table to the $T$-side factor, and by a three-coset table ($\varepsilon$,
$H$, $HS$) to the Clifford group on the $H$-side.  The two
\aid{PackedCosetTable}s over $M$ form an \aid{AmalDataNF}, and the
amalgamation module produces the alternating normal form: with $C$ the
one proper $T$-coset and $D = \{H, HS\}$, an element of
$\aid{amalNFC}\;C\;D$ is an optional leading $H$-syllable, a list of
alternations, and an optional trailing $T$ --- precisely the
Matsumoto--Amano form $(\varepsilon \mid T)\,(HT \mid SHT)^{*}\,\mathcal{C}$
read group-theoretically, as observed for this group in the Bian--Selinger
line~\cite{bian2023cliffordt2,bian2023thesis}.  The final theorem is the
syntactic isomorphism assembled by \aid{StarIsomorphism} from
generator-level translations both ways (with \aid{{\_}==={\_}} the
gate-set relations, \aid{amalPres} the amalgamated presentation, and
\aid{f} the translation):

\begin{code}%
\>[0]\AgdaFunction{clifford+T-qubit-isomorphism}\AgdaSpace{}%
\AgdaSymbol{:}\AgdaSpace{}%
\AgdaRecord{IsMonoidIsomorphism}\<%
\\
\>[0][@{}l@{\AgdaIndent{0}}]%
\>[2]\AgdaSymbol{(}\AgdaFunction{rawMonoid}\AgdaSpace{}%
\AgdaSymbol{(}\AgdaFunction{•-ε-monoid}\AgdaSpace{}%
\AgdaOperator{\AgdaDatatype{\AgdaUnderscore{}===\AgdaUnderscore{}}}\AgdaSymbol{))}\AgdaSpace{}%
\AgdaSymbol{(}\AgdaFunction{rawMonoid}\AgdaSpace{}%
\AgdaSymbol{(}\AgdaFunction{•-ε-monoid}\AgdaSpace{}%
\AgdaFunction{amalPres}\AgdaSymbol{))}\AgdaSpace{}%
\AgdaSymbol{(}\AgdaFunction{f}\AgdaSpace{}%
\AgdaOperator{\AgdaFunction{ʷ}}\AgdaSymbol{)}\<%
\end{code}
\begin{code}[hide]%
\>[0]\AgdaFunction{clifford+T-qubit-isomorphism}\AgdaSpace{}%
\AgdaSymbol{=}\AgdaSpace{}%
\AgdaFunction{CliffordT1.CliffordT1.CliffordT1-isomorphism}\<%
\end{code}

Every coset-table obligation in the tower --- well-definedness on all
axioms at all cosets, several dozen cases --- is discharged by
\aid{by-equal-nf Eq.refl}: the typechecker computes both coset
descriptors and observes that they are equal.

\subsection{Single-qutrit Clifford+$T$}

The qutrit analogue replaces $\omega$ by a ninth root $\zeta$.  The
gate-set relations (generators $T, H, S, \zeta$; abbreviations
$X := \zeta^3 (HSH)(HSSH)$ and $Z := \zeta^3 S^2 H^2 S H^2$), derived and
verified in this development, are
\[
  \zeta^9 = \varepsilon,\;\;
  S^3 = \zeta^6,\;\;
  H^4 = \varepsilon,\;\;
  T^3 = Z,\;\;
  (SH)^3 = \varepsilon,\;\;
  (THH)^2 = \varepsilon,
\]
\[
  HHSHHS = SHHSHH,\;\;
  TS = ST,\;\;
  TX = \zeta^3 S^2 X\, T,\;\;
  \zeta\ \text{central},
\]
and the verified structure is again an amalgamated free product,
\[
  \mathrm{Clifford}{+}T_3 \;\cong\;
  \langle T, M \rangle \,\ast_{M}\, \langle H, M \rangle,
  \qquad M = \langle S, X, \zeta, H^2 \rangle,
\]
reached by a longer tower: from $\langle \zeta \rangle$ through
$\langle S, \zeta \rangle$ (three cosets) and a \emph{nine}-coset table
to $\langle S, X, \zeta \rangle$ (with the twisted commutation
$(XS)(SX) = \zeta^6 (SX)(XS)$), then a two-coset table adjoining $H^2$ to
reach $M$, and finally the two extensions over $M$.  The common subgroup
$M$ is the subgroup $\mathcal{S}\mathcal{M}$ of generalized permutation
matrices (with the central phase) in the canonical form of Glaudell,
Ross, and Taylor~\cite{glaudell2019qutrit} (obtained independently by
Prakash, Jain, Kapur, and Seth~\cite{prakash2018qutrit}, and extended to
qudits in~\cite{jain2020qudit}), by which every single-qutrit
Clifford+$T$ operator is uniquely
$(\varepsilon \mid T \mid H^2T)\,
 (HT \mid H^3T \mid SHT \mid SH^3T \mid S^2HT \mid S^2H^3T)^{*}\,
 \mathcal{C}$, $T$-optimally.  Our alternating carrier is this form read
group-theoretically: the $T$-side contributes the proper cosets $T$ and
$TH^2$ (their leading factor), the $H$-side the cosets $H$, $HS$, $HS^2$
(their Clifford prefixes, in the mirrored reading), and the $2 \times 3$
alternation letters are their six syllables.  The final theorem,
\aid{clifford+T-qutrit-isomorphism} in the index, has the same shape as
the qubit one and the same proof economy: the nine-coset level alone
imposes seventy-two axiom-preservation obligations, every one closed by
\aid{refl} after evaluation.  We know of no prior machine-checked
completeness result for a qutrit gate-set presentation.

\subsection{\texorpdfstring{$U_3(\mathbb{Z}[\tfrac12,i])$}{U3(Z[1/2,i])}, a two-level amalgamation}

The third case study formalizes the case $n = 3$ of Bian and Selinger's
presentation of the groups $U_n(\mathbb{Z}[\tfrac12,i])$~\cite{bian2021un}
of unitary matrices over $\mathbb{Z}[\tfrac12,i]$.  Generators are $i_0$
(a local phase), $K_{01}$ (a scaled $2{\times}2$ Hadamard-type block),
and the transpositions $X_{01}, X_{12}$; derived generators $i_1, i_2,
K_{02}, K_{12}$ are conjugates.  The eleven relations include the orders
$i_0^4 = X_{01}^2 = X_{12}^2 = \varepsilon$, braid and commutation laws,
and the interactions of $K_{01}$ with the phases, e.g.\
$K_{01} i_1^2 = X_{01} K_{01}$ and $K_{01}^2\, i_0\, i_1 = \varepsilon$.
The verified structure is
\[
  U_3(\mathbb{Z}[\tfrac12,i]) \;\cong\;
  (\mathbb{Z}_4 \wr S_3) \,\ast_{M}\, \langle M, K_{01} \rangle,
  \qquad M = (\mathbb{Z}_4 \wr S_2) \times \mathbb{Z}_4 :
\]
the left factor is the subgroup of \emph{monomial} matrices --- the
generalized symmetric group $G(4,1,3)$, taken verbatim from the wreath
development above --- and the right factor adjoins $K_{01}$ to the common
subgroup $M$ of monomial matrices supported on wires $0, 1$ times a phase
on wire $2$; the right factor is itself presented by a six-coset
Reidemeister--Schreier level over $\langle i_0 \rangle \times \langle
i_1 \rangle$ (the tower's second level, whence \emph{two-level}), with
the derived generators introduced by definitional extensions.  Each
factor has index $3$ over $M$, with proper coset representatives
$\{X_{12},\, X_{12}X_{01}\}$ and $\{K_{01},\, K_{01} i_0\}$, so every
element is uniquely an optional $K$-syllable, a list of alternations
from their product, an optional trailing $X$-syllable, and a tail in $M$
carried by the normal form of $M$ itself.  The final theorem,
\aid{U₃-presentation-isomorphism}, is again a monoid isomorphism between
the flat presentation and the amalgamated one.  At $2\,868$ and $1\,394$
source lines, the qutrit and $U_3$ studies are the largest verified
objects in the library, and both consist mostly of first-order tables
and \aid{refl}-discharged obligations --- the shape of proof this
framework was designed to make routine.
